\chapter{Competition Events}
\newpage

\section{General rules}
A competition field consists of several concentric circles. The inner circle with a 2 metre radius is the bull’s-eye from which throws are made. The boomerang has to fly over the 20 metre circle with every throw. The number of circles used depends on the competition and there are, of course, specific rules which can differ from nation to nation. International rules are being developed.

\section{Accuracy}
In this event, the boomerang is thrown five times in a row from the centre of concentric circles with a radius of 2, 4, 6, 8, 10 and 20 metres. The boomerang must neither be caught nor touched. Points are awarded based on how close the boomerang lands to the centre. Scoring is the same as the accuracy portion of ``Australian Round'' (see \ref{fig:Scores}).

If the boomerang lands outside the 10 metre circle or the thrower touches the boomerang, the thrower receives no points. If the boomerang touches a circle line, the mid-score of the two circles is given.
For this competition, I would recommend a boomerang which flies neither too high nor too far. It should not have too much spin when it returns or it will tend to bounce usually away from the bullseye. You should use only as much strength as necessary when throwing.
\newpage

\section{Australian Round}
In Australian Round, points are awarded for accuracy, catching and distance. Catching within the 20 metre circle earns 4 points and outside earns 2 points. Points for accuracy are counted even if the boomerang isn’t caught. In addition to the circles for Accuracy there are 30, 40 and 50 metre lines. Distance points are awarded based on the largest circle the boomerang crosses before returning. Distance points are awarded only if you received points either in accuracy or in catching (see \ref{fig:Scores}).
All players make 5 throws, and the results are added together. Up to 100 points can be scored. 

When the wind is not turbulent, throw the largest range boomerang you can get bulleyes with. In windy or turbulent conditions, when accurate returns are elusive, use a 50 metre boomerang. If you catch it within the 20 meter circle you can earn a total of 50 points for range and catching, and if you catch within the 10 metre circle, you will receive additional accuracy points.

\begin{figure}[H]
\centering
\includegraphics[width=30em]{74_ARKR.eps} 
\caption{Accuracy and distance points}
\label{fig:Scores}
\end{figure}


\section{Trick Catch}
This game consists of 2 parts. 
\subsection{Part 1}
10 different catches are required. You receive points for each successful catch. The catches are as follows:
\begin{enumerate}
\item Left-hand clean catch (2 points)
\item Right-hand clean catch (2 points)
\item Two-hand catch behind the back (3 points)
\item Two-hand catch under the leg (3 points)
\item Eagle catch (4 points)
\item Hackey catch (6 points)
\item Tunnel catch (6 points)
\item One-hand catch behind the back (7 points)
\item One-hand catch under the leg (7 points)
\item Foot catch (10 points)
\end{enumerate}
It is best to use a boomerang which flies relatively high and which hovers down slowly to the thrower. Use a boomerang that has reduced spin at the end of its flight to make catching easier. If it has too little spin, however, it will be unstable and unpredictable, and more difficult to catch.
\newpage

\subsection{Part 2}
The aim is to throw two boomerangs simultaneously (figure \ref{fig:DoubleGrib}) and catch both with the following trick catches:
\begin{enumerate}
\item Catch behind the back; Catch under the leg
\item Left hand catch; Hackey catch
\item Right hand catch; Tunnel catch
\item One-hand behind the back; One-hand under the leg
\item Eagle catch; Foot Catch
\end{enumerate}

\begin{figure}[H]
\centering
\includegraphics[width=20em]{77_DOUBLE.eps} 
\caption{Double Grip}
\label{fig:DoubleGrib}
\end{figure}
\newpage

Use two boomerangs with the same point of return, but a different flight path and duration. One boomerang should return quickly, the other should hover longer. This will give you enough time to catch the second after catching the first.
If the boomerangs touch each other shortly after the launch, they will be difficult to catch (at best) ; one or both may stop spinning and fall. In order to avoid this, you can grip the ``Outsider'' by arm 1 and the ``Insider'' by arm 2 or vice versa. Or try making the Outsider heavier or swap the Outsider with the Insider (see also section \ref{sec:AchieveCorrect}). The most unusual pair of Doublers I have seen in practice, was a Fast Catch and an MTA.



\newpage
\section{Fast Catch}
The object is to throw one boomerang and catch it 5 times as quickly as possible. Time starts at the first throw and ends after the 5th catch when the thrower touches the 2 metre circle. All throws are made from within the 2 metre circle. The boomerang must pass the 20 metre line on each throw for the catch to count. The best boomerang for this competition is one which returns fast that you can catch reliably. Suitable boomerangs have many different shapes and sizes. But they all have one thing in common: they don’t have much lay-over. This can be achieved through a specific shape, (see Chapter \ref{sec:ConstructionDisk}), a certain airfoil or tuning - usually, a combination.
Fast Catch boomerangs tend to fit into two categories. One has quite a high drag and needs a strong throw - when these return, they have nearly no spin so they are easy to catch. The other has a low drag and doesn’t require a strong throw but still flies fast. Fast Catch boomerangs are often thrown with a tilt of 0° or less and on the whole, can’t take much wind.

I prefer the type of Fast Catch boomerang with a high drag. This means that the boomerang has to be thrown very hard but is harmless in the final stage of flight and easy to catch. This is an advantage for saving time. The Fast Catch should be tuned to 22 -- 23 metres. Even if the flight time is longer, you can ensure that the boomerang flies 20 metres if there is a gust of wind or if the throw is poor.
\newpage

\section{Long Distance}
Long Distance boomerangs are thrown from the centre point of a 40 metre base line. The returning boomerang must fly over this line (see figure \ref{fig:LDLine}), otherwise the throw doesn’t count. It is not necessary to catch the boomerang. Generally speaking, the heavier the model, the better it is. Long-distance boomerangs are usually weighted with lead or made out of a heavy material. The hook shape is most preferable.

\begin{figure}[H]
\centering
\includegraphics[width=30em]{75_LDOK.eps} 
\caption{Valid and unvalid flight}
\label{fig:LDLine}
\end{figure}

The difficulty in construction is that you need to produce just enough lift to support the heavy boomerang long enough for it to return. Too much lift causes a premature turn - and shorter range. You normally throw at a tilt of 45° or more and with a lot of force. Most Long Distance boomerangs need a medium to strong wind. They should only be thrown on VERY large fields. Make sure there are no people around as Long Distance boomerangs can easily get out of control! They can fly a long distance in any direction.
\newpage

\section{MTA 100 - Maximum Time Aloft in 50 metre circle}
The object of this event is to keep the boomerang in the air as long as possible, before catching it. The longest time in five tries is your score. Throw from within the 50 metre circle (100 metres in diameter, which is why the event is called MTA 100) and then catch - also within the 50 metre circle. The throw doesn’t count if thrown or caught outside this circle. Arm 1 of the conventional MTA shape is much longer than Arm 2. The inventor of this MTA shape is Wilhelm Bretfeld from Norderstedt, Germany. Ted Bailey (USA) has refined Bretfeld’s designs. Most of the MTA’s currently used in competition are based on these refinements. Ideal weather conditions for MTA are low temperatures with strong sunshine and a light wind (thermal). On such a clear day, the boomerang might even stay in the air for minutes.

If the wind isn’t too strong, you can use a wooden MTA, however, with more than 2 or 3 beauforts it is better to use a phenolic MTA. You should make the most of the 50 metre circle. The ideal position for launching is to stand on the circle line with the wind blowing from you toward the centre point. Try to get a good average time first and then risk aiming for a better time.
\newpage

\section{Juggling}
Juggling is done with two boomerangs. One boomerang should always be in the air, while the other is caught and thrown. The event continues until one is dropped. Score is the number of catches made consecutively. It is preferable to use two similar boomerangs. They should fly high and descent slowly providing plenty of time to decide where to catch and where to throw the next boomerang.





